\section{Experiments}\label{sec:experiments}
The experiments follow our plan: C1 measures the evaluation error (RQ1), C2 the solvers and grids (RQ2), C3--C5 the training-side changes (RQ3), and C6 combines them (RQ4). Each part states what was run, then what was found; every number comes from a result file named in the \LaTeX\ source and is regenerated by \texttt{report/make\_tables.py}.

\subsection{C1: How reliable is the evaluation, and what does FID miss? (RQ1)}\label{sec:c1}
\paragraph{What we ran.} FID at 1, 2, 5, 10, 20, 50 and 100 Euler (or DDIM) steps and with dopri5 for the three seeds of I-CFM and OT-CFM and the two DDPM seeds; on the seed-0 models, five further noise draws at 2, 10 and 100 steps and 2,000--20,000 samples at 100 steps; the class-coverage test at 1, 2, 5, 20 and 100 steps for every model, with three noise draws at 100 steps.
\begin{figure}[t]
\centering
\begin{minipage}[b]{0.56\linewidth}\centering
\includegraphics[width=\linewidth]{figs/fig_c1_reliability.pdf}
\end{minipage}\hfill
\begin{minipage}[b]{0.4\linewidth}\centering
\includegraphics[width=\linewidth]{figs/fig_samples.pdf}
\end{minipage}
\caption{\emph{Left (C1):} three training seeds per model. (a) FID against Euler steps (10,000 samples). (b) Total variation between generated and real $k$-means cluster shares ($k=10$); dash-dotted: held-out real test set. Lower is better in both panels; I-CFM seed~2 (dotted) has much worse FID but the same class balance as seed~0. \emph{Right:} samples from the same eight noise vectors after 1, 2, 5 and 100 steps. One I-CFM step lands on the data mean, as \eqref{eq:postmean} predicts; OT-CFM already separates digits; DDIM needs 5 steps; after reflow one step gives the digit that 100 base steps give.}
\label{fig:c1}\label{fig:samples}
\end{figure}
\begin{table}[t]
\centering
\caption{Evaluation error on MNIST (FID, lower is better). Top: base models, mean $\pm$ sd over 3 training seeds (DDIM: the 2 seeds' values), 10,000 samples. Middle: sd over 6 sampling-noise draws of the seed-0 model. Bottom: seed-0 model at NFE~100 against the number of generated samples, and the $c/n$ extrapolation FID$_\infty$ \citep{chong2020fid}.}
\label{tab:c1}
\small
\begin{tabular}{lcccc}
\toprule
 & NFE 1 & NFE 5 & NFE 100 & dopri5 \\
\midrule
I-CFM (3 seeds) & $322\pm9.28$ & $45.6\pm7.79$ & $17.0\pm5.81$ & $14.3\pm5.05$ \\
OT-CFM (3 seeds) & $212\pm9.92$ & $33.4\pm5.80$ & $16.0\pm1.88$ & $13.4\pm1.70$ \\
DDIM (2 seeds) & 293, 356 & 99.3, 106 & 12.2, 32.8 & -- \\
\midrule
sampling-noise sd & NFE 2 & NFE 10 & NFE 100 & \\
\midrule
I-CFM / OT-CFM & 0.34 / 0.58 & 0.13 / 0.28 & 0.13 / 0.05 & \\
\midrule
samples ($n$) & 2k / 5k & 10k & 20k & FID$_\infty$ \\
\midrule
I-CFM, NFE 100 & 17.3 / 14.9 & 14.5 & 14.1 & 13.7 \\
OT-CFM, NFE 100 & 17.1 / 14.5 & 13.9 & 13.4 & 13.0 \\
\bottomrule
\end{tabular}
\end{table}
% sources: results/e4/mnist_eval.json, results/e3/fid_s{1,2}.json, results/c1_reliability/fid_noise.json

\paragraph{Training seeds dominate the error.} At NFE~100 the seed sd of FID is 5.8 (I-CFM) and 1.9 (OT-CFM), against a noise sd of 0.13 and 0.05 (Table~\ref{tab:c1}, Figure~\ref{fig:c1}a), so subtracting $\sigma^2_{\text{noise}}$ changes nothing. I-CFM seed~2 is worse at every NFE (23.7 at NFE~100 versus 14.5 and 12.9) although its setup and final loss (0.169) match the other seeds; without EMA the last-epoch weights are a noisy choice. The thresholds for separately trained models are therefore large, 11.6 (I-CFM) and 3.8 (OT-CFM) at NFE~100, while samplers on one model can be compared to within a fraction of a point. With 10,000 samples FID exceeds the extrapolated FID$_\infty$ by 0.8--0.9 for both models, so comparisons at equal $n$ are unaffected; 2,000 samples add 3.6--4.0 and must not be mixed with 10,000-sample numbers. By the rule of Section~\ref{sec:eval}, OT-CFM beats I-CFM at 1, 2 and 5 steps on every seed (paired differences $-109$ to $-110$, $-49$ to $-56$ and $-1$ to $-19$) and cannot be distinguished from it from 10 steps on. Straightness is very stable: $R=1.021$--1.022 for I-CFM and 1.010--1.011 for OT-CFM. % results/c1_reliability/summary.md, fid_noise.json
\paragraph{The course method sees what FID misses.} The two-sample $\chi^2$ test does not reject the held-out real test set (TV $=0.017$, $p=0.14$, close to the floor 0.014 of \eqref{eq:tv}), whereas the one-sample test that treats $p_k$ as known wrongly rejects it ($p=0.037$). Generated samples reach TV 0.05--0.10 at 100 steps, with bootstrap intervals about 0.02 wide and a seed sd of 0.02, so TV differences below about 0.05 are not meaningful. Three findings are invisible to FID (Figure~\ref{fig:c1}b). (1) One-step I-CFM puts almost all samples into one cluster (TV 0.70--0.79; the SVM labels every sample as the same digit), as Proposition~\ref{prop:onestep} predicts, while one-step OT-CFM covers all ten digits (TV 0.08--0.14, digit shares 0.015--0.37). (2) I-CFM seed~2 has the worst FID but the same TV as seed~0 (0.098 versus 0.092): its problem is image quality, not class balance. (3) For DDIM the two criteria disagree outright: seed~0 has the best 100-step FID of all models (12.2) but the worst class balance (TV 0.19; digit 1 makes up 18.7\% of its samples, digit 3 only 4.5\%), while seed~1 is balanced (TV 0.055--0.066) at FID 32.8. Judging by FID alone, or by one seed, misleads in both directions. Another $k$-means initialisation changes TV by less than 0.006 (except one-step I-CFM, up to 0.033); with $k=20$ even the real test set is rejected (TV 0.046, $p\approx10^{-11}$), probably because finer clusters pick up style differences between the MNIST training and test sets, so $k=20$ serves only as a check. % results/e3/kmeans.json, results/c1_reliability/kmeans_c1.json

\subsection{C2: Solvers and time grids at a fixed NFE (RQ2)}\label{sec:c2}
\paragraph{What we ran.} On the seed-0 models: $\kappa(t)$; Euler, Heun and RK4 at about 2--32 NFE on the uniform, $\sqrt\kappa$, calibrated and four fixed shift grids ($s=0.25,0.5,2,4$), all with FID; endpoint errors against dopri5 for all 24 shift grids. After the first results we added the $t=1$ diagnostic (RMS of $v_\theta(x,1)-x$ on 2,000 real and 2,000 generated images), Heun* and RK4*, and repeated the main comparisons on seeds 1 and 2.
\begin{figure}[t]
\centering
\includegraphics[width=0.64\linewidth]{figs/fig_c2_solvers.pdf}
\caption{C2. (a) Path curvature $\kappa(t)=\E\|\tfrac{d}{dt}v(x(t),t)\|$ from 256 dopri5 trajectories (log scale) and the 8-step $\sqrt\kappa$ grid of OT-CFM (dotted lines); $G$ is the predicted error ratio \eqref{eq:sqrtk}. (b) FID (10,000 samples, lower is better) against NFE for OT-CFM seed~0: three solvers on the uniform (dashed) and $\sqrt\kappa$ (solid) grids, and Heun*/RK4* (crosses). Dash-dotted: dopri5.}
\label{fig:c2}
\end{figure}
\paragraph{Curvature and the $\sqrt\kappa$ grid.} For both models $\kappa(t)$ rises about a hundredfold for $t>0.95$, and 65--69\% of $\int\kappa$ lies in $t>0.9$ (Figure~\ref{fig:c2}a): the network settles fine detail last. Proposition~\ref{prop:grid} predicts error ratios $G=0.50$ (OT-CFM) and 0.54 (I-CFM). With Euler the $\sqrt\kappa$ grid lowers FID at 16 and 32 steps from 17.4 to 14.3 and from 16.2 to 12.5 (OT-CFM), and from 22.4 to 19.8 and from 18.4 to 15.2 (I-CFM); at 32 steps OT-CFM is within 10\% of dopri5 (11.5) and I-CFM within 25\% (12.1), at about 30 instead of 270 NFE. At 2--4 steps it is equal or slightly worse (Table~\ref{tab:sampler}), because it leaves one long step through the early, still-averaged part of the path (the failure mode noted in Section~\ref{sec:grid}). These are paired comparisons, so a 3-point difference is 10--60 sampling-noise sd. % results/c2_solvers/{otcfm,icfm}_s0.json
\begin{table}[t]
\centering
\caption{Inference-side samplers on the seed-0 base models: FID (10,000 samples, lower is better) at about 8, 16 and 32 network calls (exact count in parentheses when it differs), and NFE@20. The last-step variants never evaluate the network at $t=1$. Paired comparisons on one model: the sampling-noise sd is 0.05--0.6 (Table~\ref{tab:c1}).}
\label{tab:sampler}
\small
\setlength{\tabcolsep}{4pt}
\resizebox{\linewidth}{!}{%
\begin{tabular}{llcccc}
\toprule
Model & Sampler & NFE $\approx$ 8 & 16 & 32 & NFE@20 \\
\midrule
I-CFM & Euler, uniform & 32.0 & 22.4 & 18.4 & 23.7 \\
 & Euler, $\sqrt\kappa$ & 31.3 & 19.8 & 15.2 & 15.8 \\
 & Heun, uniform & 126 & 78.4 & 48.0 & $>32$ \\
 & Heun, $\sqrt\kappa$ & 61.9 & 22.1 & 13.7 & 19.0 \\
 & RK4, uniform & 125 & 65.9 & 37.1 & $>32$ \\
 & RK4, $\sqrt\kappa$ & 114 & 31.3 & 13.2 & 24.6 \\
 & Heun, last step Euler, uniform & 23.9$^{(7)}$ & 17.8$^{(15)}$ & 16.5$^{(31)}$ & 11.4 \\
 & Heun, last step Euler, $\sqrt\kappa$ & 23.8$^{(7)}$ & 16.0$^{(15)}$ & 12.9$^{(31)}$ & 10.2 \\
 & RK4, last step Heun, uniform & -- & 17.5$^{(14)}$ & 16.3$^{(30)}$ & 9.91 \\
\midrule
OT-CFM & Euler, uniform & 20.5 & 17.4 & 16.2 & 9.15 \\
 & Euler, $\sqrt\kappa$ & 19.7 & 14.3 & 12.5 & 7.88 \\
 & Heun, uniform & 134 & 83.3 & 50.4 & $>32$ \\
 & Heun, $\sqrt\kappa$ & 61.1 & 22.8 & 13.3 & 19.6 \\
 & RK4, uniform & 114 & 66.1 & 40.7 & $>32$ \\
 & RK4, $\sqrt\kappa$ & 97.0 & 31.8 & 12.8 & 24.6 \\
 & Heun, last step Euler, uniform & 18.4$^{(7)}$ & 15.5$^{(15)}$ & 15.9$^{(31)}$ & 6.60 \\
 & Heun, last step Euler, $\sqrt\kappa$ & 19.0$^{(7)}$ & 15.6$^{(15)}$ & 12.4$^{(31)}$ & 6.67 \\
 & RK4, last step Heun, uniform & -- & 16.4$^{(14)}$ & 16.5$^{(30)}$ & 7.73 \\
\bottomrule
\end{tabular}}
\end{table}
% sources: results/e4/mnist_eval.json, results/c2_solvers/{icfm,otcfm}_s0.json, results/c2_solvers/t1_{icfm,otcfm}_s0.json

\paragraph{Second-order solvers fail because of $t=1$.} On the uniform grid Heun and RK4 are far \emph{worse} than Euler at every NFE (OT-CFM at 32 NFE: 50.4 and 40.7 versus 16.2), contrary to their higher order. The diagnostic confirms the explanation of Section~\ref{sec:t1}: the networks miss the ideal value $v(x,1)=x$ by 43\% (RMS relative error) on real images and by about 85\% on generated ones, and the error of the field against its target grows from 0.36 at $t=0.9$ to 0.57 at $t=0.999$. Replacing only the last step removes the failure: at 31 and 30 NFE Heun* and RK4* reach 15.9 and 16.5 on OT-CFM instead of 50.4 and 40.7, and 16.5 and 16.3 on I-CFM instead of 48.0 and 37.1, paired differences more than 50 times the sampling noise. The best inference-only samplers reduce NFE@20 from 9.2 to 6.6 (OT-CFM, Heun*) and from 23.7 to 9.9 (I-CFM, RK4*), at no training cost. % results/c2_solvers/t1_{otcfm,icfm}_s0.json
\paragraph{Endpoint error does not predict FID.} The calibrated grid always selects steps concentrated near the noise ($s\approx0.2$--0.4) and makes FID \emph{worse} than the uniform grid in all ten Euler cases (e.g.\ 76.0 to 108.8 at two steps for OT-CFM); the endpoint-error and FID rankings of pairs of grids agree in only 124 of 204 cases (61\%). Endpoint error is dominated by early errors that decide which digit is drawn, whereas FID responds to fine detail settled near $t=1$, where $\kappa$ peaks; ODE accuracy measured without data is a poor proxy for sample quality.
\paragraph{Other training seeds.} Every direction found on seed~0 holds on all six models: the $\sqrt\kappa$ grid lowers Euler FID at 16 steps by 0.5--3.1 and at 32 steps by 1.6--3.7 and is slightly worse at 4 steps; on the uniform grid Heun is 1.9--3.7 times worse than Euler at 32 NFE, and Heun* removes that loss (14.4--24.5); the $t=1$ error is 34--55\% on real and 83--91\% on generated images. The size of the $\sqrt\kappa$ gain varies between models, so we report it per model.

\subsection{C3--C5: Training-side changes and their cost (RQ3)}\label{sec:c35}
\paragraph{What we ran.} C3: the pairing cost saved by exact OT at OT batches from 16 to 4,096 (CPU), and OT-CFM trained with OT batches 512 and 2,048. C4: I-CFM trained with logit-normal and $\kappa$-weighted times. C5: reflow from the seed-0 I-CFM and OT-CFM models. Each variant is compared with its base by the threshold (Table~\ref{tab:training}; bold entries exceed it).
\begin{table}[t]
\centering
\caption{Training-side variants, training seed 0, Euler sampler on the uniform grid, FID on 10,000 samples (lower is better). $R$: path-length ratio (1 = straight). NFE@20: steps to reach FID 20. Time: training wall time (reflow: plus generating the pairs) relative to the base model, both throttled to the same duty cycle (OT batch 2,048 was resumed after a power loss at epoch 20; its time adds the 20 epochs before the resume). Bold: differs from the base model of the same method by more than the improvement threshold $2\,\mathrm{sd}$ over training seeds (last row).}
\label{tab:training}
\small
\setlength{\tabcolsep}{2.9pt}
\begin{tabular}{lccccccccc}
\toprule
 & \multicolumn{5}{c}{FID at NFE} & & & & \\
\cmidrule(lr){2-6}
Model & 1 & 2 & 5 & 10 & 100 & dopri5 & $R$ & NFE@20 & Time \\
\midrule
I-CFM (base) & 332 & 132 & 45.3 & 27.9 & 14.5 & 12.1 & 1.021 & 23.7 & 1.00 \\
OT-CFM (base) & 223 & 76.0 & 26.7 & 19.0 & 13.9 & 11.5 & 1.011 & 9.15 & 1.00 \\
OT-CFM, OT batch 512 & \textbf{195} & 70.4 & 30.6 & 23.5 & 17.4 & \textbf{16.7} & 1.010 & 22.1 & 1.30 \\
OT-CFM, OT batch 2,048 & 208 & 82.9 & 37.8 & 26.4 & 17.4 & \textbf{15.6} & 1.009 & 31.9 & 1.97 \\
I-CFM, logit-normal $t$ & 327 & \textbf{103} & 30.0 & 18.3 & 12.6 & 14.2 & 1.020 & 9.06 & 1.22 \\
I-CFM, $\kappa$-weighted $t$ & 335 & 122 & 37.2 & 26.0 & 18.3 & 17.5 & 1.024 & 58.1 & 1.10 \\
I-CFM + reflow & \textbf{64.6} & \textbf{30.9} & \textbf{23.9} & 22.6 & 20.5 & 19.0 & 1.001 & $>100$ & 2.67 \\
OT-CFM + reflow & \textbf{61.2} & \textbf{28.9} & 23.8 & 22.6 & \textbf{21.5} & \textbf{21.1} & 1.001 & $>100$ & 2.83 \\
\midrule
\multicolumn{10}{l}{Threshold $2\,\mathrm{sd}$ over three training seeds, I-CFM / OT-CFM:} \\
 & 18.6/19.8 & 10.2/7.81 & 15.6/11.6 & 14.1/7.58 & 11.6/3.76 & 10.1/3.41 & & & \\
\bottomrule
\end{tabular}
\end{table}
% sources: results/e4/mnist_eval.json, results/e3/fid_s0_*.json, results/e2/*_train.json, results/c5_reflow/pairs_*.json

\paragraph{C3: Minibatch OT saves little at any batch size.} Exact OT saves 2.9\% of the random pairing cost at batch 16, 4.5\% at 128, 5.2\% at 512, 5.8\% at 2,048 and 6.1\% at 4,096, i.e.\ 11--24\% of the bound \eqref{eq:bound}. Each doubling saves about 0.4\% more; extrapolating this logarithmic trend, 10\% would need batches of about $2.8\times10^6$, beyond the 60,000 training images, while one OT solve already takes 1.2\,s at batch 2,048 and 6.2\,s at 4,096 (\texttt{results/c3\_otbatch/pairs.json}). Training confirms that this is too little to matter. OT batch 512 has a one-step FID of 196, lower than all three base seeds (205--223), but the gain exceeds the threshold only against seed~0, the worst base seed at one step, not against the three-seed mean (16.3 versus 19.8); from five steps on it is 3--5 points worse. OT batch 2,048 has the straightest OT-CFM paths ($R=1.0085$) and the lowest loss (0.100 versus 0.118, a smaller constant $C$ in Proposition~\ref{prop:cfm}), yet its one-step FID (208) lies inside the base-seed range and it is 3--11 points worse than seed~0 from two steps on; it doubled the training time, 20\% of which was OT solving. % results/e3/fid_s0_otb{512,2048}.json
\paragraph{C4: Where to train matters more than how to pair.} Logit-normal times, which rarely sample $t$ near 0 or 1, give FID 103, 30.0, 18.3, 14.1 and 12.0 at 2, 5, 10, 20 and 50 steps, better than all three base seeds at every NFE from 2 to 100, and lower NFE@20 from 23.7 to 9.1, as much as the best sampler on the base model (9.9). By the strict rule only the two-step gain (29.4, threshold 10.2) is detectable; at five steps the gain (15.3) is just below the threshold (15.6). A fourth draw from the seed distribution would rank first at any one NFE with probability $1/4$, but the ranks at different NFEs are strongly correlated, so this is suggestive rather than a test, and it rests on one training seed. Consistent with less training near $t=1$, where the logit-normal density vanishes, its dopri5 FID (14.2, at only 111 NFE), which evaluates the network at $t=1$, is worse than its own 50-step Euler FID (12.0). Drawing half of the times where $\kappa$ is large gives the curviest paths of all models ($R=1.024$; dopri5 needs 345 NFE instead of 280) and no detectable change in FID: the sampler should spend its steps near $t=1$, but training should not spend its samples there. % results/e3/fid_s0_t{logitnormal,kappa}.json
\paragraph{C5: Reflow straightens the paths but lowers the ceiling.} We integrated each seed-0 model with dopri5 from 60,000 fresh noise vectors (287 and 277 NFE per pair) and fine-tuned it on these pairs. The pairs cost 1,319 and 1,316, a saving of 12.7--12.9\% of the random cost and half of the bound \eqref{eq:bound}, twice what minibatch OT achieves at batch 4,096. The loss falls from 0.168 to 0.004, as Proposition~\ref{prop:cfm} predicts for a deterministic coupling ($C\approx0$). The paths become almost straight ($R\approx1.001$), and few-step FID improves far beyond the threshold: one step from 332 to 64.6 (I-CFM) and from 223 to 61.2 (OT-CFM), two steps from 132 to 30.9 and from 76.0 to 28.9. With many steps, however, both models plateau near 21 (dopri5: 19.0 and 21.1 instead of 12.1 and 11.5), beyond the threshold for OT-CFM from 20 steps on. This is the ceiling of Section~\ref{sec:reflow}: the reflowed model can at best reproduce its base model's ODE samples and adds the error of a second regression; we have not separated this from the effect of fine-tuning rather than training from scratch. After reflow I-CFM and OT-CFM are almost identical (differences of at most 3.4), so OT-CFM's base advantage of about 110 one-step FID points does not survive reflow. Class balance is unharmed: I-CFM reflow has cluster TV 0.066--0.069 at every NFE (base model: 0.79, 0.25 and 0.13 at 1, 2 and 5 steps, 0.085--0.092 at 100). Reflow costs 2.7 and 2.8 times the base training time (pair generation 2.05 and 1.84\,h plus retraining 1.08 and 1.56\,h, against 1.17 and 1.20\,h). % results/e3/fid_s0_reflow.json, results/c5_reflow/pairs_*_s0.json, kmeans_reflow.json

\subsection{C6: Combining the two sides (RQ4)}\label{sec:c6}
\paragraph{What we ran.} The core C2 study (Euler, Heun and RK4 on the uniform and $\sqrt\kappa$ grids) on both reflow models, and the straightness ratio of all twelve models.
\begin{figure}[t]
\centering
\includegraphics[width=0.66\linewidth]{figs/fig_frontier.pdf}
\caption{C6, quality--efficiency frontier, training seed~0 (FID, 10,000 samples, lower is better). Solid/dash-dotted: Euler on the uniform grid; dotted: improved samplers on the $\sqrt\kappa$ grid (Heun* for base models, Euler for reflow). Grey: lower envelope. Variants without a detectable gain are omitted.}
\label{fig:frontier}
\end{figure}
\paragraph{Samplers stack with reflow only partly.} On both reflow models the $\sqrt\kappa$ grid lowers Euler FID from 4 steps on by 0.5--2.4 (I-CFM reflow: 21.7 to 19.4 at 32 steps; OT-CFM reflow: 24.7 to 22.3 at 4). On the uniform grid Heun and RK4 still fail although the paths are nearly straight (Heun at 32 NFE: 46.2 and 34.6, against 21.7 and 21.9 for Euler), so the failure is not caused by curvature. We did not run the $t=1$ diagnostic on the reflow models; for a deterministic coupling the ideal $u_1(x)=x-\E[x_0\mid x_1=x]$ is no longer $x$. % results/c2_solvers/{icfm,otcfm}_s0_reflow.json
\paragraph{The best change depends on the budget.} The lower envelope of Figure~\ref{fig:frontier} is formed by reflow at 1--4 NFE (61.2, 28.9 and, with the $\sqrt\kappa$ grid, 22.3 at 4) and, from 7 NFE on, alternately by OT-CFM with Heun* or Euler on the $\sqrt\kappa$ grid and by the logit-normal I-CFM model with plain Euler (19.0 at 7, 14.1 at 20, 12.4 at 31, 12.0 at 50); these curves are within about one FID point of each other, well below the seed threshold, so the envelope does not rank them. Reflow buys a 2.6--3.6 times lower one- and two-step FID than the best base model for 2.7--2.8 times the training time but never reaches FID 20; the samplers cut NFE@20 by 28--58\% for free. The same pattern appears in the published CIFAR-10 results of Table~\ref{tab:literature}.
\paragraph{Straightness separates groups, not models.} Across the twelve models the Spearman correlation between $R$ and one-step FID is 0.92 (0.85 without the reflow models), driven by the gaps between groups (reflow $<$ OT-CFM $<$ I-CFM); within OT-CFM and within I-CFM (five models each) it is only 0.3 and 0.5. The models share seed~0, so these numbers are descriptive.
